\documentclass[conference]{IEEEtran}
\usepackage[T1]{fontenc}
\usepackage[english,brazil]{babel}
\usepackage{amsmath,amssymb,booktabs,graphicx,cite,url}
\graphicspath{{../results/figures/}}
\newcommand{\imagunit}{\mathrm{j}}
\newcommand{\figfile}[2]{\includegraphics[width=#1]{#2}}
\author{%
\IEEEauthorblockN{Gabriel Gon\c{c}alves Brand\~ao, Maria Roberta Alves Brand\~ao,\\
Lucas Estanislau da Costa e Kennedy Fernandes Dos Santos}
\IEEEauthorblockA{Pontif\'{i}cia Universidade Cat\'olica de Goi\'as -- Escola Polit\'ecnica e de Artes}
\InputIfFileExists{matriculas.local.tex}{}{}
}

% Generated by paper/build_tables.py from results/processed/.
\newcommand{\ChecksPass}{129}
\newcommand{\ChecksTotal}{129}
\newcommand{\DelayCritical}{$1{,}20920$}
\newcommand{\NominalRows}{A & $(0,0,0)$ & $10{,}000$ & $180{,}00$ \\
B & $(1,-1,0)$ & $0{,}500$ & $19{,}44$ \\
C & $(0,0,0)$ & $3{,}000$ & -- \\
D & $(0,0,0)$ & $1{,}322$ & $40{,}61$ \\}
\newcommand{\DelayRows}{1 & $2{,}0000000$ & $65{,}39867\%$ & estável \\
2 & $1{,}2360680$ & $2{,}22200\%$ & estável \\
3 & $1{,}2100529$ & $0{,}07057\%$ & instável \\
5 & $1{,}2091998$ & $0{,}00001\%$ & instável \\}
\newcommand{\CpuRows}{Nyquist & $1{,}007$ & $1{,}903$ & $1{,}621$ \\
Routh--Hurwitz & $0{,}071$ & $0{,}231$ & $0{,}307$ \\
Lugar das raízes & $6{,}708$ & $16{,}459$ & $23{,}054$ \\
Bode & $0{,}789$ & $1{,}580$ & $2{,}424$ \\}

\title{Análise Computacional e Comparativa de Critérios de Estabilidade em Sistemas Dinâmicos}

\begin{document}
\maketitle
\renewcommand{\abstractname}{Resumo}
\begin{abstract}
Este trabalho implementa e compara Nyquist, Routh--Hurwitz, lugar das raízes
e Bode para laços lineares, contínuos e de uma entrada e uma saída. O núcleo
em Python avalia a resposta complexa, conta envoltórios em um contorno fechado
e calcula cruzamentos e margens. O estudo usa quatro plantas que representam
um modelo estável de terceira ordem, uma planta instável, um zero no semiplano
direito e um atraso de transporte. Medições em ordens 5, 10 e 20 distinguem
tempo de execução, condicionamento e erro em raízes. A comparação com soluções
analíticas e ferramentas independentes permite verificar a contagem de polos,
as margens e a aproximação de Padé. Os \ChecksPass{} controles de validação passaram. A fronteira de
atraso do caso D foi de \DelayCritical{} s. Padé de ordens 1 e 2
ainda indicou estabilidade a 1\% acima dessa fronteira.

\end{abstract}
\renewcommand{\IEEEkeywordsname}{Palavras-chave}
\begin{IEEEkeywords}
estabilidade, Nyquist, Routh--Hurwitz, lugar das raízes, Bode, atraso de transporte.
\end{IEEEkeywords}
{\selectlanguage{english}
\renewcommand{\abstractname}{Abstract}
\begin{abstract}
This work implements and compares Nyquist, Routh--Hurwitz, root locus, and
Bode analysis for continuous-time, single-input single-output feedback loops.
A Python core evaluates the complex response, counts encirclements on a
closed contour, and computes crossover frequencies and stability margins.
Four plants cover an open-loop stable third-order model, an unstable plant,
a right-half-plane zero, and a transport delay. Experiments at orders 5, 10,
and 20 distinguish runtime, conditioning, and root error. Analytical
solutions and independent software provide reference values for pole counts,
margins, and Padé approximation. All \ChecksPass{} validation checks passed. The case D delay boundary was
1.20920 s. First- and second-order Padé models still indicated stability at
1\% above that boundary.

\end{abstract}
\renewcommand{\IEEEkeywordsname}{Keywords}
\begin{IEEEkeywords}
stability, Nyquist, Routh--Hurwitz, root locus, Bode, transport delay.
\end{IEEEkeywords}
}
\section{Introdução}
Um sistema de realimentação pode estabilizar uma planta com um polo no
semiplano direito. O mesmo laço pode perder estabilidade quando a comunicação
introduz atraso suficiente para alterar a fase relativa entre entrada e
saída. Esse efeito aparece em sistemas de controle em rede, nos quais o
tempo de transporte, a amostragem e a variação temporal da comunicação
afetam a malha \cite{zhang2001}. A pergunta investigada é em que condições
quatro métodos clássicos concordam sobre estabilidade e que erros emergem
quando a análise exige amostragem de frequências ou operações com
polinômios de ordem elevada.

A comparação usa uma implementação que mantém a análise de controle no
núcleo e isola as referências computacionais na validação. Os experimentos
avaliam o efeito de um polo instável, de um zero de fase não mínima e do
atraso puro. A ferramenta expõe os valores usados no diagnóstico e oferece
uma interface para variar os parâmetros e reproduzir os diagramas.

\section{Trabalhos relacionados}
Nyquist formulou a análise de regeneração em função da resposta de malha
\cite{nyquist1932}. Bode sistematizou o desenho de amplificadores realimentados
por magnitude e fase \cite{bode1945}. A contagem algébrica de Routh
\cite{routh1877} e o método geométrico de Evans \cite{evans1950} abordam a
mesma questão a partir do polinômio característico. Ogata e Nise organizam
esses métodos em uma sequência para sistemas contínuos
\cite{ogata2010,nise2017}. O capítulo de análise em frequência de Åström e
Murray relaciona Nyquist, margens e sensibilidade \cite{astrom2008}.
Para o atraso, a aproximação racional de Padé dá um caminho para métodos
polinomiais \cite{baker1996}. Em ordem elevada, a sensibilidade das raízes
a perturbações nos coeficientes pode dominar o erro observado
\cite{wilkinson1959}. Esses resultados orientam a escolha das referências
e das métricas deste estudo.

\section{Fundamentação teórica}
Considere a realimentação unitária negativa de uma planta contínua própria,
com coeficientes reais, ganho escalar $K$ e atraso $\tau\ge 0$.
\begin{equation}
L(s)=K\frac{B(s)}{A(s)}e^{-s\tau},\qquad
T(s)=\frac{L(s)}{1+L(s)} .
\label{eq:loop}
\end{equation}
Na parte racional, os polos e zeros são as raízes de $A$ e $B$,
respectivamente. Com $\tau=0$, o polinômio característico é
$C(s)=A(s)+KB(s)$. A análise preserva a ordem de $A$ e os fatores comuns,
para que um modo interno cancelado na análise entrada--saída permaneça
no modelo algébrico. Para estabilidade assintótica, os polos de malha
fechada devem ficar no semiplano esquerdo.

\subsection{Mapeamento e princípio do argumento}
O contorno segue o eixo de $-\imagunit R$ a $+\imagunit R$ e volta pelo arco à direita.
Indentam-se polos sobre o eixo com semicírculos que os excluem da
região contada. Para $F=1+L$, a variação de argumento no contorno
horário é $-2\pi(Z-P)$, onde $P$ e $Z$ contam polos e zeros no interior.
O mapeamento por $L$ escreve a mesma contagem em torno de $-1$.
\begin{equation}
\begin{aligned}
Z &= N+P,\\
N &= -\frac{1}{2\pi}\sum_{k}\arg
 \left(\overline{u_k}\,u_{k+1}\right),\\
u_k &= 1+L(s_k).
\end{aligned}
\label{eq:nyquist}
\end{equation}
A fórmula discreta fecha o polígono e preserva a convenção de que
envoltórios horários contribuam positivamente. Ela aproxima a
variação contínua quando o mapeamento entre amostras é resolvido
\cite{astrom2008}.

\subsection{Routh--Hurwitz, lugar das raízes e Bode}
A tabela de Routh dispõe os coeficientes alternados nas duas primeiras
linhas. Se $r_{i,j}$ é a entrada da linha $i$ e coluna $j$,
\[
r_{i,j}=r_{i-2,j+1}
-\frac{r_{i-2,0}}{r_{i-1,0}}\,r_{i-1,j+1}.
\]
As mudanças de sinal na primeira coluna contam raízes no semiplano
direito. Um pivô nulo exige tratar o limite de um valor positivo.
Uma linha inteira nula exige a derivada do polinômio auxiliar
\cite{routh1877,nise2017}. O lugar das raízes mostra o conjunto das
raízes de $A+KB$ para ganhos sucessivos e permite observar
amortecimento, frequência natural e cruzamentos \cite{evans1950}.

O diagrama de Bode mostra $20\log_{10}|L(\imagunit\omega)|$ e a fase
desenrolada. Em um cruzamento de fase $\omega_p$ no semieixo real
negativo, a margem de ganho é $1/|L(\imagunit\omega_p)|$. Em um cruzamento
de módulo unitário $\omega_g$, a margem de fase é o afastamento
angular de $-180^{\circ}$ na convenção local. Um ganho com margem
inferior a um pode indicar uma fronteira ao reduzir $K$. Quando
existem vários cruzamentos, cada um descreve uma fronteira candidata.
A classificação global usa a contagem de Nyquist ou os polos, conforme
o modelo \cite{bode1945,astrom2008}.

\subsection{Atraso e aproximação racional}
A exponencial acrescenta $-\omega\tau$ rad à fase e preserva o módulo
no eixo. A fronteira pode ocorrer em $\omega>0$ quando
$|L_0(\imagunit\omega)|=1$ e
\begin{equation}
\arg L_0(\imagunit\omega)-\omega\tau=(2q+1)\pi,
\qquad q\in\mathbb{Z}.
\label{eq:delay}
\end{equation}
Aqui $L_0$ é a função sem atraso. Para aproximar o atraso no plano
$s$, o Padé diagonal de ordem $m$ usa $Q_m(-s\tau)/Q_m(s\tau)$.
Se $c_k$ são os coeficientes crescentes de $Q_m$, então
$c_0=1$ e
\[
c_{k+1}=c_k\frac{m-k}{(k+1)(2m-k)}.
\]
A série coincide com a da exponencial até a potência $2m$.
O módulo dessa aproximação no eixo é unitário. A fase e os polos
fechados dependem da ordem e da banda de avaliação \cite{baker1996}.

\section{Metodologia e arquitetura}
O pacote separa modelo de planta, métodos de análise, validação e
aplicação. O núcleo usa Python 3.12 e NumPy. O sistema define
coeficientes em ordem decrescente e avalia a razão por Horner, com
uso de polinômios reversos para $|s|>1$. A instrução de teste verifica
as importações permitidas no núcleo. SciPy, Python-control, SymPy e
Mpmath são executados na camada de validação.

\subsection{Contorno numérico e margens}
A malha contém frequências positivas e negativas, os arcos de desvio
e o arco externo. Para $A$ mônico, o código usa
$1-\sum_{i\ge1}|a_i|R^{-i}$ como cota inferior normalizada do
denominador no exterior, quando positiva. A cota superior da razão
é calculada sobre a mesma escala. Para um modelo estritamente
próprio e com atraso, $|e^{-s\tau}|\le1$ no semiplano direito permite
conter a cauda em um disco que exclui $-1$. Sem atraso, a cota se
aplica à diferença em relação ao ganho de passagem direta. Com atraso,
um ganho de passagem direta de módulo maior ou igual a um fica fora
da certificação desse contorno.

A contagem exige duas malhas sucessivas com o mesmo inteiro,
incremento angular máximo inferior a $\pi/3$ e distância do polígono
ao ponto crítico acima de $10^{-7}$. A distância é calculada nos
segmentos e nos vértices. O refinamento local bissecta segmentos cujo
mapeamento varia rapidamente ou apresenta curvatura relevante em relação
à distância ao ponto crítico. Cruzamentos de ganho unitário são pontos
adicionais da malha. A resposta próxima a $-1$ em um cruzamento
recebe o estado de fronteira. Um resultado sem resolução suficiente
recebe diagnóstico indeterminado. Quando a distância de um polo aberto ao eixo é numericamente
ambígua, a contagem fica indeterminada. Os raios de indentação
são reduzidos para preservar polos fechados próximos ao eixo.
Em sistemas racionais, a distância às raízes de $A+KB$ fornece
essa escala. Em sistemas atrasados, verifica-se a dominância do
termo de realimentação sobre o denominador em torno do polo.
A cota externa limita a região de interesse. A resolução no
interior continua empírica.

Para o ganho unitário, calcula-se o polinômio em $x=\omega^2$
a partir de $|KB(\imagunit\omega)|^2-|A(\imagunit\omega)|^2$. Cada candidato
encontrado é verificado na resposta complexa. Para a fase racional,
o produto $KB(\imagunit\omega)\overline{A(\imagunit\omega)}$ identifica a
interseção com o eixo real. Com atraso, a busca por semieixo negativo
ocorre em uma banda finita explícita. As passagens de Bode são
registradas em rad/s e as margens em razão de ganho e graus.

\subsection{Métodos polinomiais e resposta ao degrau}
Routh calcula a recorrência em ponto flutuante. Pivôs pequenos e
linhas nulas acionam uma alternativa racional para reduzir cancelamento.
A contagem no eixo usa as raízes nulas e os fatores comuns das partes
par e ímpar do polinômio. O lugar das raízes calcula as raízes
do polinômio para cada $K$ e conecta as amostras por atribuição
de custo mínimo sobre distâncias no plano complexo. O cálculo de
cruzamentos verifica os candidatos do produto de polinômios em
$\imagunit\omega$. Em caso de queda da ordem, a trajetória pode possuir
um ramo no infinito. O degrau usa uma realização companheira e uma
exponencial matricial calculada por série e quadraturas sucessivas.
Para plantas atrasadas, o degrau, Routh e lugar das raízes usam a
mesma aproximação de Padé escolhida pelo usuário.

\subsection{Plantas e hipóteses}
A Tabela~\ref{tab:systems} resume as plantas. Em A, os polos são
$-1,-2,-3$ e a fronteira de ganho é $K=10$ em $\sqrt{11}$ rad/s.
Em B, os polos são $1$ e $-2$. O modelo representa uma dinâmica
linearizada com rigidez efetiva negativa, como um pêndulo
invertido amortecido em torno do equilíbrio vertical. A escrita
normalizada é $\ddot x+\dot x-2x=u$. A fronteira em $K=2$ é
um cruzamento por frequência zero. Em C, os polos são $-1,-2$
e o zero é $+1$. O degrau inicia com inclinação negativa.
O polinômio $s^2+(3-K)s+(2+K)$ prevê cruzamento em $K=3$ e
$\sqrt5$ rad/s. Em D, o polo racional é $-1$ e a
primeira fronteira de atraso no ganho nominal é
$\tau_*=2\pi/(3\sqrt3)$ s, obtida de \eqref{eq:delay}.
As plantas A, B e D não têm zeros finitos.

\begin{table}[!t]
\centering
\caption{Plantas de teste e parâmetros nominais.}
\label{tab:systems}
\begin{tabular}{clcc}
\toprule
Caso & $B(s)/A(s)$ & $K$ & $\tau$ (s)\\
\midrule
A & $6/(s^3+6s^2+11s+6)$ & 1 & 0\\
B & $1/(s^2+s-2)$ & 4 & 0\\
C & $(1-s)/(s^2+3s+2)$ & 1 & 0\\
D & $2/(s+1)$ & 1 & 0,8\\
\bottomrule
\end{tabular}
\end{table}

\subsection{Validação e medições}
Os comparadores usam resultados analíticos onde há uma fronteira
conhecida e métodos de bibliotecas fora do núcleo. Assim é possível
registrar diferenças entre a representação em coeficientes e a avaliação
da função, a contagem de raízes e as convenções sobre cruzamento em DC.
Para os benchmarks, CPU e tempo de parede são medidos após aquecimento.
A carga de módulos, a geração de gráficos e a escrita em disco ficam
fora da seção cronometrada. Dados brutos guardam cada repetição e a
descrição do ambiente. A memória é medida à parte, sobre alocações
Python. A sensibilidade em raízes simples é interpretada pela relação
de primeira ordem $|\delta r|\simeq |\delta C(r)|/|C'(r)|$,
com distinção entre erro frente à construção analítica e resíduo
escalado do polinômio armazenado \cite{wilkinson1959}.

\section{Resultados e validação}
\subsection{Estabilidade absoluta e relativa}
Foram aprovadas \ChecksPass{} de \ChecksTotal{} verificações
computacionais nos experimentos. A Tabela~\ref{tab:margins}
apresenta os pontos nominais. Routh também classificou os quatro
casos nominais como estáveis, usando Padé $[3/3]$ para D.
A comparação incluiu configurações abaixo, sobre e acima das
fronteiras. Em A, o ganho $K=10$ levou um par ao eixo. Em B,
$K=2$ produziu um polo na origem. Em C, o par imaginário
surgiu em $K=3$. O cálculo independente do lugar das raízes
concordou com esses três pontos de ganho. As
Figs.~\ref{fig:nyquist}--\ref{fig:locus} mostram os diagramas
nos quatro sistemas. A Fig.~\ref{fig:margins} compara as margens
nominais na convenção da tabela.

\begin{table}[!t]
\centering
\caption{Estabilidade e margem nos pontos nominais.}
\label{tab:margins}
\begin{tabular}{cccc}
\toprule
Caso & $(P,N,Z)$ & MG (razão) & MF ($^\circ$)\\
\midrule
\NominalRows
\bottomrule
\end{tabular}
\end{table}

No caso A, a margem de fase da tabela decorre de um cruzamento
em DC. O comparador Python-control reporta a margem de ganho
de 10, mas sua rotina de margens usuais omite esse cruzamento
unitário em DC. A diferença foi registrada nas convenções
experimentais. Em B, a margem de ganho de $0{,}5$ identifica a
fronteira ao reduzir o ganho nominal de 4 para 2. A contagem
$P=1$ e $N=-1$ demonstra a estabilidade em malha fechada.
Em C, inexiste um cruzamento de ganho unitário na frequência
real e a MG é 3. A resposta inicial inversa produz um mínimo
amostrado de $-0{,}19487$, embora a amostra final se aproxime
de $1/3$ no intervalo avaliado. A Fig.~\ref{fig:step} mostra a
resposta ao degrau dos quatro casos.

O erro relativo máximo da resposta complexa nominal frente à
referência foi de $6{,}22\times10^{-16}$. O maior erro absoluto
nos conjuntos de raízes por ganho foi de $3{,}56\times10^{-14}$
e a maior diferença no degrau foi de $7{,}25\times10^{-14}$
nas quatro plantas, dentro das tolerâncias do protocolo.
Para D, a referência temporal representa o mesmo modelo racional com Padé $[3/3]$. O acordo nesse degrau mede o algoritmo de resposta
temporal e não mede erro físico do Padé.

\begin{figure*}[!t]
\centering
\figfile{0.95\textwidth}{nyquist-abcd.pdf}
\caption{Contornos de Nyquist das plantas A--D nos pontos nominais.
A cruz assinala $-1$. O caso B exige um envoltório anti-horário.}
\label{fig:nyquist}
\end{figure*}

\begin{figure*}[!t]
\centering
\figfile{0.97\textwidth}{bode-abcd.pdf}
\caption{Magnitude e fase de A--D. O caso D usa o atraso exponencial.
A faixa exibida de frequências é finita.}
\label{fig:bode}
\end{figure*}

\begin{figure*}[!t]
\centering
\figfile{0.95\textwidth}{root-locus-abcd.pdf}
\caption{Lugar das raízes em A--D. Em D, os ramos são os
do modelo racional de Padé $[3/3]$.}
\label{fig:locus}
\end{figure*}

\begin{figure*}[!t]
\centering
\figfile{0.96\textwidth}{stability-margins.pdf}
\caption{Margens dos quatro sistemas nominais. Para D,
selecionou-se o primeiro cruzamento de fase positivo em frequência.}
\label{fig:margins}
\end{figure*}

\begin{figure}[!t]
\centering
\figfile{\columnwidth}{step-response.pdf}
\caption{Degrau unitário de malha fechada. A curva C apresenta
deslocamento inicial negativo. Para D usa-se Padé $[3/3]$.}
\label{fig:step}
\end{figure}

\subsection{Atraso exato e Padé}
Em D, o menor cruzamento de módulo unitário positivo fica em
$\sqrt3$ rad/s. O cálculo encontrou a fronteira em
\DelayCritical{} s, em concordância com \eqref{eq:delay}.
No ponto nominal, a MF é $40{,}61^\circ$. O primeiro cruzamento
de fase positivo ocorre em $2{,}44822$ rad/s e fornece a MG
linear de $1{,}32229$. Em frequências mais altas existem
cruzamentos adicionais, com margens correspondentes ao módulo
cada vez menor da planta.

Em $\tau=5$ s, a primeira versão com refinamento global
encontrou quatro envoltórios, mas a variação de argumento
entre amostras excedeu seu critério e o resultado ficou
indeterminado. A bissecção localizada resolveu a transição
com a configuração padrão. Contaram-se $P=0$, $N=4$ e
$Z=4$, confirmados por dois pares à direita na solução por
Lambert $W$. A execução arquivou o diagnóstico original e
os valores antes e depois do refinamento.

A Tabela~\ref{tab:pade} compara a primeira fronteira dos modelos
racionais com o atraso exato. Em $1{,}01\tau_*$, as ordens
1 e 2 ainda indicam estabilidade, embora o contorno exato
conclua instabilidade. As ordens 3 e 5 recuperam a classificação
nesse ponto. O erro complexo e a discrepância de classificação
podem ser observados na Fig.~\ref{fig:pade}. A ordem adequada
depende da banda e da distância da fronteira.

\begin{table}[!t]
\centering
\caption{Atraso crítico estimado por Padé para D.}
\label{tab:pade}
\begin{tabular}{cccc}
\toprule
Ordem & $\hat\tau_*$ (s) & Erro relativo & Em $1{,}01\tau_*$\\
\midrule
\DelayRows
\bottomrule
\end{tabular}
\end{table}

\begin{figure*}[!tp]
\centering
\figfile{0.88\textwidth}{delay-pade.pdf}
\caption{Diferenças de frequência e de atraso crítico no caso D.
A referência usa a exponencial.}
\label{fig:pade}
\end{figure*}

\subsection{Desempenho e erro numérico}
As medições foram executadas em macOS 26.6.2, em um Apple M5 com
10 núcleos lógicos e 16 GiB. Usaram-se dois aquecimentos e
11 repetições por ordem e método, num total de 132 medições.
A varredura do lugar das raízes contém 48 ganhos e a de Bode,
1001 frequências. A semente da perturbação dos coeficientes
é 24112026. As amostras, os valores de CPU e parede, a média,
a mediana e o desvio-padrão estão preservados para auditoria.

\begin{table}[!t]
\centering
\caption{Mediana do tempo de CPU em milissegundos.}
\label{tab:cpu}
\begin{tabular}{lrrr}
\toprule
Método & $n=5$ & $n=10$ & $n=20$\\
\midrule
\CpuRows
\bottomrule
\end{tabular}
\end{table}

A Tabela~\ref{tab:cpu} e a Fig.~\ref{fig:bench} registram o crescimento
observado no protocolo. A carga de Routh é uma tabela por
polinômio. Lugar das raízes calcula conjuntos de raízes e
associa trajetórias para 48 ganhos. Nyquist avalia múltiplos
contornos e conta envoltórios. Bode calcula margens e amostra
duas curvas. Essas diferenças de tarefa condicionam a interpretação
das medições. Os valores de memória usam \texttt{tracemalloc}
em execução separada, refletindo o pico de alocações rastreadas
pelo Python.

\begin{figure*}[!tp]
\centering
\figfile{0.88\textwidth}{scalability.pdf}
\caption{Escalabilidade no protocolo de 11 repetições por combinação.
A dispersão foi medida com gráficos construídos dos dados brutos.}
\label{fig:bench}
\end{figure*}

O experimento de precisão comparou duas famílias de raízes
construídas, uma dispersa e uma agrupada. Em ordem 20, os erros
absolutos máximos em relação às raízes de projeto foram
$0{,}13280$ e $0{,}31928$ s$^{-1}$. Os máximos de
condicionamento relativo foram $7{,}59\times10^{14}$ e
$2{,}27\times10^{24}$. As duas famílias exibiram resíduos
escalados inferiores a $3\times10^{-16}$. Esses valores mostram
que uma raiz pode satisfazer numericamente o polinômio com
resíduo pequeno e ainda estar distante da raiz de projeto.
A referência em precisão estendida separa o arredondamento
dos coeficientes do erro do solucionador no polinômio
armazenado. A experiência repetiu o cálculo após perturbação
relativa de $10^{-10}$ nos coeficientes.
Na Fig.~\ref{fig:precision}, a dispersão dos resultados depende
da distribuição das raízes e da ordem considerada.

\begin{figure*}[!tp]
\centering
\figfile{0.88\textwidth}{numerical-precision.pdf}
\caption{Erro de localização e condicionamento para as famílias
dispersa e agrupada em três ordens.}
\label{fig:precision}
\end{figure*}

\section{Aplicação prática}
A interface em Streamlit e Plotly permite variar em conjunto
os coeficientes ou os polos e zeros, o ganho e o atraso.
Nyquist, Bode, Routh, lugar das raízes e resposta ao degrau
consomem as mesmas entradas e mostram a origem do modelo
racional ou exponencial usado. A exportação JSON conserva
a classificação, a convenção empregada, as margens e as amostras.
Testes de integração exercitam a troca de plantas, a atualização
dos dados e a condição de realimentação algébrica singular.

Uma bancada simulada separada usa a planta $1/(s+1)$ e um PID
com derivada filtrada, de parâmetros $K_p=2$, $K_i=1$,
$K_d=0{,}1$ e $T_f=0{,}04$ s. O modelo contínuo da função de
laço é
\begin{equation}
L_{\mathrm{PID}}(s)=
\frac{0{,}18s^2+2{,}04s+1}{0{,}04s^3+1{,}04s^2+s}e^{-s\tau}.
\end{equation}
A primeira fronteira prevista por Nyquist foi $1{,}05906$ s a
$1{,}78789$ rad/s. Em 90 s simulados a 50 Hz, os atrasos
constantes quantizados de 0,80 e 1,32 s produziram, respectivamente,
acomodação e crescimento do erro. Os RMS finais foram
$1{,}16\times10^{-9}$ e aproximadamente $3{,}76\times10^3$.
A repetição a 100 Hz preservou essa distinção. Nos cenários
com variação uniforme de até $\pm0{,}10$ s, o desvio-padrão
do atraso na fila foi próximo de 0,058 s. O comportamento
observado na simulação segue na Fig.~\ref{fig:hil}. A previsão
de Nyquist se refere ao atraso constante, enquanto a simulação
com jitter inclui discretização, fila e descarte de comandos
fora de ordem. Os resultados correspondem a um modelo linear
sem saturação. O código de comunicação para bancada permanece
sem medição em dispositivo físico.

\begin{figure*}[!t]
\centering
\figfile{0.89\textwidth}{hil-pid.pdf}
\caption{Erro no PID simulado para atrasos constantes e variáveis,
abaixo e acima da fronteira contínua prevista.}
\label{fig:hil}
\end{figure*}

\section{Discussão e limitações}
O estudo confirma que os métodos têm domínios e custos distintos.
Routh responde com uma contagem absoluta em um polinômio e usa
$O(n^2)$ operações de tabela na recorrência usual. O custo em
aritmética racional depende do crescimento de numeradores e
denominadores. Lugar das raízes resolve a dinâmica em vários
ganhos. Com matrizes densas e atribuição cúbica, o modelo de
custo típico é $O(gn^3)$ para $g$ amostras de ganho. O erro
pode crescer perto de raízes repetidas e a associação depende
da resolução da varredura.

A avaliação por Horner em $M$ pontos custa $O(Mn)$ e a
contagem de envoltórios custa $O(M)$ para uma malha dada.
O número de pontos depende do refinamento, da escala dos
polos, do atraso e da distância a $-1$. O cálculo dos polos
abertos e candidatos de cruzamento adiciona custo de raízes
polinomiais. Em Bode, a busca de margens racionais também
resolve polinômios em $\omega^2$, e no caso atrasado precisa
delimitar a banda de cruzamentos de fase. Essas diferenças
impedem atribuir superioridade por uma única medida de tempo.

Os resultados concentram-se em sistemas contínuos SISO e
em uma estação de trabalho. A validação exclui atrasos
neutros fora da cota adotada, cancelamentos numericamente
ambíguos e garantias de estabilidade sob jitter arbitrário.
Uma malha refinada pode perder uma transição entre amostras.
A estimativa de precisão para polinômios agrupados indica
que escala, tolerância e precisão de coeficientes merecem
atenção mesmo quando o estado de estabilidade concorda.
Em uma aplicação com dispositivo físico, seria necessário
medir atraso, jitter e latência de processamento, incluir
saturação e comparar séries em hardware.

\section{Conclusão}
A classificação nominal foi consistente nos quatro sistemas.
A contagem em B tornou explícito como um laço aberto
instável pode ser estabilizado pelo ganho. O zero no
semiplano direito de C apareceu no comportamento transitório,
e a comparação em D quantificou o erro das aproximações de
Padé perto da fronteira. A concordância dos métodos depende
das hipóteses do modelo, da banda e da precisão numérica
disponível. A ferramenta relaciona o diagnóstico aos
contornos, margens e raízes que o sustentam e disponibiliza
os dados para repetir a análise. Ensaios físicos e métodos
de contorno com cotas locais mais fortes são extensões do
trabalho.

\section*{Declaração de uso de inteligência artificial}
O OpenAI Codex foi usado para propor a metodologia, implementar e
revisar o código, executar testes e experimentos, e redigir e revisar
o resumo e as seções do artigo. Os gráficos foram gerados por scripts
a partir dos dados da execução. O repositório registra os comandos,
as fontes e os resultados de referência para conferência.
\IEEEtriggeratref{6}
\bibliographystyle{IEEEtran}
\bibliography{references}
\end{document}
